% ------------------------------------
% theorem 样式
% ------------------------------------
\declaretheoremstyle[
spaceabove=6pt, spacebelow=6pt,
headfont=\normalfont\bfseries,
headpunct=,
postheadspace=0.5em,
bodyfont=\normalfont,
headformat=\NAME
]{mycolored}

\declaretheorem[style=mycolored, name=, within=section]{theorem}
\declaretheorem[style=mycolored, name=, sibling=theorem]{lemma}
\declaretheorem[style=mycolored, name=, sibling=theorem]{corollary}
\declaretheorem[style=mycolored, name=, sibling=theorem]{proposition}
\declaretheorem[style=mycolored, name=, sibling=theorem]{definition}
\declaretheorem[style=mycolored, name=, sibling=theorem]{remark}
\declaretheorem[style=mycolored, name=, sibling=theorem]{example}
\declaretheorem[style=mycolored, name=, sibling=theorem]{construction}
\declaretheorem[style=mycolored, name=, sibling=theorem]{problem}
\declaretheorem[style=mycolored, name=, sibling=theorem]{axiom}
\declaretheorem[style=mycolored, name=, sibling=theorem]{prp}
\declaretheorem[style=mycolored, name=, sibling=theorem]{note}
\declaretheorem[style=mycolored, name=, sibling=theorem]{ill}
\declaretheorem[style=mycolored, name=, sibling=theorem]{algo}

% ------------------------------------
% 修正编号：title 中使用计数器+1
% ------------------------------------
\newcommand{\NextTheoremNum}{\thesection.\number\numexpr\value{theorem}+1\relax}
\newcommand{\NextLemmaNum}{\thesection.\number\numexpr\value{lemma}+1\relax}
\newcommand{\NextCorNum}{\thesection.\number\numexpr\value{corollary}+1\relax}
\newcommand{\NextPropNum}{\thesection.\number\numexpr\value{proposition}+1\relax}
\newcommand{\NextDefNum}{\thesection.\number\numexpr\value{definition}+1\relax}
\newcommand{\NextRemNum}{\thesection.\number\numexpr\value{remark}+1\relax}
\newcommand{\NextExNum}{\thesection.\number\numexpr\value{example}+1\relax}
\newcommand{\NextConNum}{\thesection.\number\numexpr\value{construction}+1\relax}
\newcommand{\NextProblemNum}{\thesection.\number\numexpr\value{problem}+1\relax}
\newcommand{\NextAxiomNum}{\thesection.\number\numexpr\value{axiom}+1\relax}
\newcommand{\NextPrpNum}{\thesection.\number\numexpr\value{prp}+1\relax}
\newcommand{\NextNoteNum}{\thesection.\number\numexpr\value{note}+1\relax}
\newcommand{\NextIllNum}{\thesection.\number\numexpr\value{ill}+1\relax}
\newcommand{\NextAlgoNum}{\thesection.\number\numexpr\value{algo}+1\relax}

% ------------------------------------
% 彩色 tcolorbox 定理环境
% 内部白色：colback=white，只用边框颜色
% 仅蓝 / 绿 / 橙三色
% ------------------------------------
% 蓝色系：定理、推论、命题、公理、空真原理
\AtBeginEnvironment{theorem}{%
	\theoremFont
		
\AtBeginEnvironment{corollary}{%
	\theoremFont
		
\AtBeginEnvironment{proposition}{%
	\theoremFont
		
\AtBeginEnvironment{axiom}{%
	\theoremFont
		

\AtBeginEnvironment{lemma}{%
	\theoremFont
		
% 橙色系：定义、注记、示例、构造、记号、说明
\AtBeginEnvironment{definition}{%
	\theoremFont
		
\AtBeginEnvironment{remark}{%
	\theoremFont
		
\AtBeginEnvironment{example}{%
	\par\theoremFont
	\noindent\textbf{示例 \NextExNum.}\quad
}
\AtEndEnvironment{example}{\par}

\AtBeginEnvironment{construction}{%
	\theoremFont
		
\AtBeginEnvironment{note}{%
	\theoremFont
		
% 问题/题目环境：AccentBlue 区分于定理
\AtBeginEnvironment{problem}{%
	\theoremFont
	
\AtBeginEnvironment{ill}{%
	\theoremFont
		
% 算法环境
\AtBeginEnvironment{algo}{%
	\theoremFont
		
% ------------------------------------
% Solution 环境：无色框，与 proof 相同，不编号
% ------------------------------------
\newenvironment{solution}{%
	\par\proofFont
	\noindent\textbf{解答:}\quad
}{%
	\hfill $\square$\par
}

% ------------------------------------
% 证明环境：无色块 + 中文楷体
% ------------------------------------
\renewenvironment{proof}[1][证明]{%
	\par\proofFont
	\noindent\textbf{#1.}\quad
}{%
	\hfill $\square$\par
}
% 添加星号版本（无编号）
\declaretheorem[numberwithin=chapter, name=定理, sibling=theorem]{theorem*}
\declaretheorem[numberwithin=chapter, name=问题, sibling=problem]{problem*}
\setstretch{1.2}


% ------------------------------------
% 其它设置
% ------------------------------------
\setstretch{1.2}
\geometry{textheight=9in,textwidth=5.5in,top=1in}
